% ============================================
% LUC-Entrevista
% Gerado de molambudos.md — build_miolo.py
% ============================================
\fragtitle{luc}{L}{LUC}{Entrevista}{A Entrevista}
\markboth{LUC-Entrevista --- A Entrevista}{}

\noindent \textit{O homem que viu o ciclo começar}

\noindent \textbf{Abril de 2026 --- Barbacena, Casa do Sr. Antenor}

\vspace{0.5em}\noindent\rule{\textwidth}{0.5pt}\vspace{0.5em}

\noindent O endereço estava no arquivo do Colônia, anotado a lápis na borda de um relatório de 1958: \textit{\textquotedblleft{}Antenor Francisco da Silva, ex-enfermeiro, Rua da Matriz 42, Barbacena.\textquotedblright{}}

\noindent Lúcia bateu na porta às 10 da manhã. Um homem muito velho abriu. Ele devia ter mais de 90 anos --- a pele fina como papel de seda, os olhos azuis desbotados, as mãos trêmulas.

\noindent \textit{\textquotedblleft{}A senhora é a psiquiatra?\textquotedblright{}}

\noindent \textit{\textquotedblleft{}Sou.\textquotedblright{}}

\noindent \textit{\textquotedblleft{}Então leu o diário.\textquotedblright{}}

\noindent Não era uma pergunta.

\noindent \textit{\textquotedblleft{}...Li.\textquotedblright{}}

\noindent O velho Antenor abriu a porta devagar: \textit{\textquotedblleft{}Então entre. Antes que o cheiro entre junto.\textquotedblright{}}

\vspace{0.5em}\noindent\rule{\textwidth}{0.5pt}\vspace{0.5em}

\noindent A casa era pequena, escura, cheirando a mofo e remédio. Na parede da sala, uma fotografia amarelada: homens de jaleco branco na frente do Hospital Colônia, década de 1950.

\noindent Antenor serviu café. A xícara tremia na mão dele.

\noindent \textit{\textquotedblleft{}O que a senhora quer saber?\textquotedblright{}}

\noindent \textit{\textquotedblleft{}Tudo o que sabe sobre o paciente 1.260.\textquotedblright{}}

\noindent \textit{\textquotedblleft{}Joaquim.\textquotedblright{}\textit{ Ele disse o nome como quem diz o nome de um parente morto. }\textquotedblleft{}Conheci ele. Conheci todos eles. Mas o Joaquim era diferente.\textquotedblright{}}

\noindent \textit{\textquotedblleft{}Diferente como?\textquotedblright{}}

\noindent \textit{\textquotedblleft{}Ele não era louco. A senhora leu o diário. Sabe que ele não era louco.\textquotedblright{}}

\noindent \textit{\textquotedblleft{}Eu sei.\textquotedblright{}}

\noindent \textit{\textquotedblleft{}Ele via coisas.\textquotedblright{}}

\noindent \textit{\textquotedblleft{}Que tipo de coisas?\textquotedblright{}}

\noindent \textit{\textquotedblleft{}A criatura. Ele via a criatura. Desde 1915. Antes de entrar no Colônia. Ele já via.\textquotedblright{}}

\noindent Antenor contou a história de Joaquim como quem conta uma história que ouviu muitas vezes. Joaquim não era violento quando chegou --- foi a vala que o tornou violento. A vala onde os corpos não enterrados eram comidos por bichos. A vala onde Joaquim viu a criatura pela primeira vez.

\noindent \textit{\textquotedblleft{}Ele me disse uma vez: 'A criatura não veio de fora. Ela já estava na vala. Enterrada junto com os corpos. Eu só... acordei ela.'\textquotedblright{}}

\noindent \textit{\textquotedblleft{}E o senhor acreditou?\textquotedblright{}}

\noindent Antenor olhou para Lúcia. Os olhos azuis desbotados, mas a lucidez absoluta.

\noindent \textit{\textquotedblleft{}Eu trabalhei 42 anos no Colônia. Eu vi pacientes terem convulsões. Eu vi pacientes serem lobotomizados. Eu vi pacientes morrerem de fome. Eu vi coisas que a senhora não acreditaria se eu contasse. Mas de todas as coisas que vi, a única que eu não consegui explicar foi o diário do Joaquim.\textquotedblright{}}

\noindent \textit{\textquotedblleft{}O que o senhor quer dizer?\textquotedblright{}}

\noindent Ele se levantou. Foi até um armário. Tirou uma caixa de sapatos. Abriu.

\noindent Ela achou que ia encontrar fotos, documentos, lembranças.

\noindent Em vez disso, encontrou um caderno. Capa preta. Páginas amareladas. Letra trêmula.

\noindent \textit{\textquotedblleft{}Isso é...\textquotedblright{}}

\noindent \textit{\textquotedblleft{}O segundo diário de Joaquim.\textquotedblright{}}

\noindent Lúcia sentiu o sangue gelar.

\noindent \textit{\textquotedblleft{}Ele escreveu dois?\textquotedblright{}}

\noindent \textit{\textquotedblleft{}A senhora leu o primeiro. O que ele escreveu vivo. Este ele escreveu morto.\textquotedblright{}}

\noindent \textit{\textquotedblleft{}Como?\textquotedblright{}}

\noindent \textit{\textquotedblleft{}Eu cavei a cova dele. Três dias depois do enterro, a cova estava remexida. Abri. O caixão estava quebrado por dentro. O caderno estava no peito dele, com a mão dele em cima. A letra era a mesma. Mas a data era depois da morte.\textquotedblright{}}

\noindent Ele entregou o caderno a Lúcia.

\noindent \textit{\textquotedblleft{}A senhora pode ficar com ele. Eu morri de medo dele a vida inteira. Já é hora de passar adiante. É assim que funciona, não é?\textquotedblright{}}

\noindent Lúcia segurou o caderno. A capa estava fria. Mas, dentro dela, algo se aqueceu.

\noindent \textit{\textquotedblleft{}É assim que funciona\textquotedblright{}}, ela repetiu.

\noindent \textit{\textquotedblleft{}Sim\textquotedblright{}\textit{, disse Antenor. }\textquotedblleft{}A senhora já sabe. A senhora só não aceitou ainda.\textquotedblright{}}

\vspace{0.5em}\noindent\rule{\textwidth}{0.5pt}\vspace{0.5em}

\vspace{0.5em}
\noindent\textit{\textbf{↪ Rotas:} \rota{LUC-Escolha} \rotasep \rota{LUC-Contraprova}}